

\section{课题来源与选题依据}

\subsection{研究背景}

\subsubsection{6G网络的智能化演进与安全挑战}

移动通信网络正经历从5G向6G的代际演进。根据 ITU-R IMT-2030 愿景~\cite{ITU_IMT2030_2023}，未来 IMT 系统将面向 2030 年及以后
支持更广泛的应用场景与服务需求。在此基础上，IMT-2030(6G) 推进组白皮书进一步提出，6G 网络将实现“万物智联”（Intelligence of Everything）~\cite{IMT2030_6G_WhitePaper_2021}，并在无线接入、核心网、边缘计算等各层深度融合人工智能（AI），形成``AI原生网络''（AI-Native Network）架构。在此架构中，AI不再仅作为辅助工具用于资源调度和网络优化，而是成为网络运行的核心驱动力，承担异常检测、威胁识别、风险评估等关键安全功能。

然而，AI在网络安全领域的深度应用引入了新的风险维度。现有研究表明~\cite{BCNN-IDS,Tabarisaadi2024}，基于深度学习的入侵检测系统（Intrusion Detection System, IDS）在训练集覆盖的攻击类型上可达到95\%以上的准确率，但在面对零日攻击（Zero-Day Attack）、变种攻击或数据分布漂移时，模型可能以高置信度输出错误预测。这种现象源于确定性神经网络的固有特性：模型将所有输入强制映射到已知类别，即使输入与训练分布存在显著差异，仍会输出单一的确定性预测，而无法表达其内在的认知不确定性（Epistemic Uncertainty）。

\subsubsection{零信任架构对可信AI的需求}

6G网络安全架构的核心理念是零信任（Zero-Trust Architecture, ZTA）~\cite{NIST-ZTA}，其基本原则为``持续验证、永不信任''（Never Trust, Always Verify）。在零信任框架下，安全决策必须满足``可解释、可度量、可追责''的要求，这对AI驱动的安全系统提出了更高的可信性标准。具体而言，可信AI（Trustworthy AI）在网络安全场景中需要具备以下核心能力：

\begin{enumerate}
  \item \textbf{不确定性感知能力}：模型应能识别其预测不确定的样本，而非对所有输入均输出表面上高置信度的单一概率分布。这要求模型能够区分``确信某样本为正常流量''与``不确定但被迫分类为正常流量''两种本质不同的情况。

  \item \textbf{风险量化能力}：模型应为每个预测提供可度量的置信度和风险评估指标，支持安全运维人员进行告警优先级排序和人工复核决策。传统确定性模型输出的softmax概率仅反映类别间的相对关系，并不能真实反映模型对预测的信心程度。

  \item \textbf{分布外检测能力}：模型应能识别与训练分布显著不同的输入（Out-of-Distribution, OOD），避免对未知攻击类型做出盲目预测。这对于应对快速演化的网络威胁至关重要。
\end{enumerate}

上述需求在传统确定性深度学习框架下难以满足，这构成了本课题研究的根本动机。

\subsubsection{网络攻击的时空演化特性}

网络攻击并非孤立的数据包事件，而是在时间和空间维度上呈现结构化模式的行为序列。从时间维度看，攻击行为表现为前后关联的演化模式：端口扫描呈现递增的端口序列，分布式拒绝服务（DDoS）攻击表现为突发的流量峰值，慢速攻击则展现长时间的低频异常行为。从空间维度看，相似的攻击数据包在特征空间中形成聚类结构：相同协议类型、相似包长度、相同传输方向的数据包往往属于同一攻击阶段。

现有入侵检测方法在建模这种时空双重依赖关系时存在局限性。基于循环神经网络（RNN）或Transformer的时序模型专注于捕捉包序列的前后关联，但难以利用相似包的空间聚类信息；基于图神经网络（GNN）的方法能够建模网络拓扑或流量关系图，但往往假设静态结构，无法充分捕捉攻击的时间演化特性。这种建模能力的不足限制了模型对复杂攻击模式的识别能力。

\subsection{研究意义}

\subsubsection{理论意义}

本课题的理论意义体现在以下三个方面：

\textbf{（1）贝叶斯深度学习与时空图神经网络的融合创新}

现有贝叶斯深度学习研究主要聚焦于计算机视觉和自然语言处理领域，在网络安全场景中的应用尚处于初步探索阶段~\cite{BCNN-IDS}。同时，时空图神经网络在入侵检测中的应用多采用确定性学习范式~\cite{E-ResGAT,TCG-IDS}，缺乏对预测不确定性的系统建模。本课题将变分推断（Variational Inference）引入时空图卷积网络，构建端到端的贝叶斯时空图神经网络框架，在理论上实现了两个研究方向的深度融合，为可信AI在网络安全领域的应用提供了新的方法论基础。

\textbf{（2）包级时空图建模范式的创新}

现有时空图方法多在主机级或流级粒度建模~\cite{EvolveGCN,PPT-GNN}，需要维护流状态表和预聚合统计特征，在高并发场景下面临计算复杂度和内存开销瓶颈。本课题提出的包级时空图建模范式，通过动态图构建和时空并行卷积，在无需流状态维护的前提下直接在数据包序列上捕捉时空依赖关系，在理论上突破了现有方法的粒度限制，为实时入侵检测提供了新的技术路径。

\textbf{（3）不确定性量化理论在开放世界安全场景中的扩展}

传统不确定性量化研究多假设封闭世界（Closed-World）设定，即测试数据与训练数据服从相同分布。然而，网络安全场景本质上是开放世界（Open-World）问题：新型攻击不断涌现，数据分布持续演化。本课题通过区分认知不确定性（Epistemic Uncertainty）和偶然不确定性（Aleatoric Uncertainty），并将其与分布外检测、主动学习等机制结合，在理论上扩展了不确定性量化方法在开放世界场景中的适用性。

\subsubsection{实践意义}

本课题的实践意义体现在以下四个方面：

\textbf{（1）提升6G网络的安全防护能力}

6G网络的超高速率（峰值速率达1 Tbps）、超低时延（空口时延<1 ms）和超大连接（每平方公里百万级设备）特性，对入侵检测系统的实时性和准确性提出了极高要求。本课题提出的BSTGCN框架通过包级时空图建模和自适应采样策略，在保持高检测精度的同时将推理延迟控制在可接受范围内（<100 ms/batch），为6G网络提供了可落地的安全解决方案。

\textbf{（2）支持零信任架构的可信决策}

零信任架构要求安全系统能够为每个访问请求提供动态的风险评估。本课题通过不确定性量化机制，使模型能够输出可度量的置信度指标，支持基于风险的访问控制决策。例如，对于高不确定性的流量，系统可触发多因素认证或人工审查，而对于低不确定性的正常流量则允许快速通过，从而在安全性和用户体验之间实现平衡。

\textbf{（3）降低安全运维成本}

现有入侵检测系统普遍存在高误报率问题，导致安全运维人员需要处理大量虚假告警。本课题通过不确定性驱动的告警优先级机制，使运维人员能够优先处理高风险、高不确定性的告警，显著提升运维效率。实验表明，基于不确定性的优先级策略可使运维人员在处理前20\%告警时即覆盖80\%的真实攻击。

\textbf{（4）增强模型的持续学习能力}

网络威胁持续演化，入侵检测模型需要具备持续学习能力以适应新型攻击。本课题通过主动学习机制，利用不确定性指标选择最具信息量的样本进行标注和再训练，在有限标注预算下最大化模型性能提升。这为构建可持续演进的安全系统提供了技术支撑。

\subsection{研究问题的提出}

基于上述背景和意义分析，本课题聚焦于以下核心研究问题：

\textbf{研究问题1：包级时空图建模}

在无需维护流状态表的前提下，如何在数据包序列上构建时空图表示，同时捕捉攻击行为的时间演化（前后包的关联）和空间聚类（相似包的结构）特性，并保持包级时间分辨率以支持实时检测？

\textbf{研究问题2：贝叶斯推断与不确定性量化}

如何在时空图神经网络中引入变分推断机制，使模型在保持高分类精度的同时输出可信的不确定性度量，并区分认知不确定性和偶然不确定性，同时兼顾实时检测场景下的计算效率要求？

\textbf{研究问题3：不确定性驱动的安全决策}

如何将不确定性量化从模型内部指标转化为可操作的安全决策支持机制，包括告警优先级排序、人工复核策略、分布外检测和主动学习等，以提升入侵检测系统的实用性和可信度？

上述三个研究问题构成了本课题的核心研究内容，其解决将为6G网络安全中的可信入侵检测提供系统性解决方案。



\subsection{相关研究与技术空缺}

\subsubsection{时空图神经网络在入侵检测中的应用}

网络攻击在时间和空间上均呈现结构化特征：攻击行为随时间演化（如多阶段渗透），同时在网络拓扑中传播（如横向移动）。近年来，图神经网络（GNN）因其对关系结构的建模能力，被广泛引入网络入侵检测（NIDS）~\cite{lo2022egraphsage,E-ResGAT,wang2025bsgat}。现有研究可按建模粒度划分为三个层次：主机级、流级和包级，各层次在时空建模能力、检测延迟和计算开销上呈现不同权衡。

\paragraph{（1）主机级时空图：捕捉网络拓扑中的攻击传播}


主机级方法将网络中的终端设备（主机、服务器）建模为图节点，通信关系建模为边，通过GNN捕捉攻击在拓扑中的传播模式。早期工作如E-GraphSAGE~\cite{lo2022egraphsage}将物联网流量表示为主机间通信图，节点表示IP地址和端口组合，边表示通信流量，边特征为流的统计属性（如平均字节长度、包计数、最小/最大/平均包到达间隔等），采用GraphSAGE聚合邻居信息进行边分类检测，在UNSW-NB15、BoT-IoT等数据集上相比传统机器学习方法平均提升2\%的加权F1分数。E-ResGAT~\cite{E-ResGAT}将网络表示转换为线图，其中流被表示为节点，应用注意力机制聚合二跳邻域并更新中心节点表示。然而，这两种方法都只考虑从大型静态网络图中随机采样的子图，未能捕捉网络流量的动态特性。BS-GAT~\cite{wang2025bsgat}进一步引入图注意力网络（GAT）构建行为相似性图，在缓解过拟合的同时将二分类F1-score提升至99\%。Sun等人~\cite{sun2024gnnids} 将领域知识图谱引入IDS，利用MulVAL生成的攻击图作为网络拓扑，将实时流量指标映射到图节点，采用GCN/GAT等图神经网络识别受攻击主机。为了提升系统的可解释性，作者采用了事后可解释性分析技术，通过识别对模型预测贡献最大的关键子图节点及其连接关系（即入侵路径），来解读模型的决策依据，从而清晰标识出完整的攻击链。

为捕捉攻击的时间演化，近期工作将GNN与时序模型结合。EvolveGCN~\cite{EvolveGCN}作为通用的动态图神经网络模型（最初发表于KDD 2020），使用门控循环单元（GRU）更新GCN权重，使模型随时间步演化，可应用于包括网络安全在内的多种动态图场景；Rossi等人提出的Temporal Graph Networks (TGN)~\cite{TGN}为连续时间动态图建模奠定了基础，该模型发表于ICML 2020，同样属于通用时序图框架。TCG-IDS~\cite{TCG-IDS}将网络流量按时间窗口构建为动态图序列（节点为终端设备，边为通信关系），采用TGN编码器并设计三重对比学习策略（时序对比、非对称对比、掩码对比）学习节点行为模式，在CSE-CIC-IDS2018上达到99.48\%平衡准确率和0.25\%误报率。在宏观攻击场景下，Euler~\cite{king2023euler}、CONTINUUM~\cite{bahar2025continuum}等系统利用时空GNN建模企业网络中的多阶段横向移动或APT攻击链，实现长时间尺度攻击行为的检测与追踪。

主机级方法的优势在于直接建模网络拓扑，适合检测分布式攻击（如DDoS、横向移动），但其时间分辨率通常在秒级，依赖预聚合的流统计特征，难以捕捉单个流内部的精细演化（如握手、重传、微突发），且在5G时代缺乏针对新兴攻击类型的验证。

\paragraph{（2）流级图嵌入：早期检测的探索与局限}


为缓解传统流特``只能在流结束后统计、导致检测滞''的问题，流级图嵌入方法将每条网络流（由五元组标识）的前缀数据包建模为图，通过无监督图嵌入或图神经网络实现早期检测。Hu等提出的Graph2vec+RF~\cite{Graph2vec+RF}代表了这一思路的早期探索：将双向流的前$n$个数据包构建为流图（节点为数据包，边为连续分组的通信突发），通过Graph2vec无监督嵌入学习图特征，再用随机森林分类，在仅利用10个初始数据包时在CIC-IDS2017上达到0.97准确率，CIC-IDS2018上达到0.993准确率。Graph2vec采用类似doc2vec的思想，将图视为文档，将图中所有节点周围的根子图视为单词，通过skip-gram模型和负采样策略学习图嵌入，专注于图的子结构特征而非节点属性，适合流图这种节点属性较少的场景。此后，一些工作在Graph2vec+RF的基础上设计更复杂的流图构造、异构图结构或自监督预训练，以提升在小样本和跨数据集场景下的表现。这些方法强调了在流前缀上构建图表示对于及时发现攻击的重要性，相比传统依赖完整流统计的方法，大大降低了检测延迟。

然而，流级图嵌入在工程实践中面临三方面关键瓶颈：

\textbf{(a) 流状态维护的计算与内存开销}——这类方法以五元组（源/目的IP、端口及传输层协议）为键，必须为每条活跃流维护独立的流状态表和图结构。系统在每个新包到达时需要：(1)查表定位其所属流，(2)更新该流对应的图，(3)当流前缀达到预定长度后计算图嵌入并交由分类器预测。随着并发流数量$M$的增加，流表查找和状态维护的开销在实际实现中近似随$O(M)$增长，并且需要在内存中同时保存成千上万条流的中间状态，这在高并发5G/6G场景下会成为严重瓶颈。

\textbf{(b) 短流攻击的上下文不足}——这类方法在实际运行中真正能够利用的上下文，本质上仍然只是``流前缀''。对于大量只包含极少数数据包的短流攻击（如 DDoS 探测报文、端口扫描、SYN Flood 等），难以在攻击早期累积足够多的数据包来构建有代表性的流级图结构，从而导致早期检测效果高度不稳定。随着数据中心与移动网络（包括 5G/6G 场景）速率的持续提升，越来越多业务可以在极短时间内完成传输，短流占比进一步攀升。已有 Internet 与数据中心流量测量普遍表明，80\% 以上的网络流由少于 10 个数据包组成，且绝大多数为低时延、小体量的 mice flows~\cite{Pekar2021}，这使得依赖长流上下文的图建模方法在短流攻击检测中面临先天不足。


\textbf{(c) 确定性模型的过度自信}——这些方法多采用随机森林或其他确定性分类器，虽可输出``概率分数''，但并未显式建模预测不确定性，对未知攻击模式仍可能过度自信。在零信任架构下，系统需要对每个预测给出可信度评估，以支持风险感知的接入控制决策，而现有流级图嵌入方法难以满足这一需求。

\paragraph{（3）包级时空图：从流内建模到跨流时空依赖}


为进一步降低检测延迟并捕捉更精细的时空模式，包级方法直接在数据包序列上构建图，无需等待流结束或累积统计特征。这一方向可分为两类：流内包级图和跨流时空图。

流内包级图方法仍以单条流为基本单元，但在流内部构建包级图结构。Ghadermazi等人（2025）的GTAE-IDS框架~\cite{GTAE-IDS}将网络流建模为动态图，用前2-8个数据包构建图来实现实时异常检测。GTAE-IDS采用无监督的图Transformer自编码器（Graph Transformer-based Autoencoder）提取流量的全局图特征，并通过重构误差检测异常。在图构建方面，GTAE-IDS将每条流的源和目的地址作为节点，数据包作为边特征，采用50维特征编码包方向（前25维表示前向包，后25维表示后向包），通过零填充区分前向和后向数据包。在CIC-IDS2018、CIC-IDS2017和ACI-IoT-2023等数据上其平均真阳性率超过98\%且误报率很低，实现了仅凭最初几个数据包即可精准识别攻击。从建模形式上看，这种方法已经接近``时空图神经网络''：节点为源/目的地址，边特征承载包级信息和方向，编码器通过多头注意力显式建模"空间依赖+时间顺序"。然而，GTAE-IDS仍然需要为每条流维护包序列并构造子图，在高并发环境中同样面临流表维护和内存开销问题；同时，其核心模型为确定性自编码器，虽适合异常检测，但尚未针对多类别攻击识别、跨场景迁移和不确定性量化展开系统设计。

跨流时空图方法则突破单流限制，在多条流或全局包序列上构建时空图。Van Langendonck等人提出的PPT-GNN~\cite{PPT-GNN}通过在NetFlow级别构建小时间窗（0.5-5秒）内的流图序列，将GNN与时间卷积/循环结构结合，在NF数据集上实现小窗口内的近实时入侵检测。PPT-GNN采用滑动时间窗口方法，为每个窗口内的所有流和IP构建空间图，并引入两种时间组件：(1)窗口内时间边（intra-window edges），将同一源或目的IP的流按时间顺序连接；(2)窗口间时间边（inter-window edges），连接跨窗口的重复节点以捕捉长期时间依赖。通过自监督预训练（负边采样和链接预测任务），PPT-GNN在UNSW-NB15、ToN-IoT和BoT-IoT数据集上相比E-ResGAT和E-GraphSAGE平均提升10.38\%的多类F1分数，并且预训练模型可以轻松迁移到新网络场景，显著降低部署成本。这些工作证明：相较于单一时间切片或纯流特征，时空图神经网络能够更好地捕捉``谁在何时与谁交互''的复杂模式，在高维、异构的网络安全场景中具有明显优势。

尽管包级方法在时间分辨率上具有理论优势，但现有工作仍存在两方面根本性不足：\textbf{(a) 大多数方法仍依赖流级抽象}——无论是流内包级图还是跨流时空图，大多数方法仍需为每条流维护包序列并构造子图，或者在流级特征上构建时间快照，时间分辨率通常在秒级。对单个流内部精细的包级演化（例如几个RTT内的握手、重传、微突发）仍然难以建模，也往往依赖预先聚合好的统计特征，仍然存在一定检测延迟。此外，现有评估方法存在不现实的问题：许多模型在大型静态图上随机采样子图进行训练和验证，这种随机子采样会破坏结构和动态关系，导致数据泄漏（来自协同攻击的流可能出现在不同的分割中），并且采用节点掩码的转导学习设置，导致模型难以泛化到新的、未见过的数据。\textbf{(b) 缺乏不确定性量化机制}——核心模型多为确定性自编码器或分类器，这些框架几乎全部采用点估计方式训练GNN，无法给出结构化的置信区间和风险指标。在面对未知攻击类型、跨时间分布漂移或对抗样本时，模型可能给出高置信度的错误预测，缺乏自我提醒能力。

\paragraph{（4）小结：时空图神经网络的核心局限}

综合上述三个层次的研究，时空图神经网络在入侵检测中已取得显著进展，但在不确定性量化方面存在系统性空缺：现有方法几乎全部采用确定性神经网络，输出的``概率''仅为softmax归一化结果，并不反映模型的真实不确定性~\cite{BCNN-IDS,Tabarisaadi2024}。这导致模型在面对未知攻击、分布漂移或对抗样本时，仍可能给出高置信度的错误预测，无法满足零信任架构对可信AI的要求。因此，如何将贝叶斯深度学习与时空图神经网络结合，成为该领域的关键技术空缺。


\subsubsection{贝叶斯深度学习在安全场景中的探索}

贝叶斯深度学习（Bayesian Deep Learning, BDL）通过在神经网络权重上引入概率分布，将点估计扩展为分布估计，从而量化模型的预测不确定性。近年来，BDL在网络安全领域的应用逐渐兴起，但大多聚焦于传统的流统计特征或图像化流量表示，尚未与时空图神经网络深度结合。

\paragraph{（1）贝叶斯神经网络在入侵检测中的初步探索}

Haider等人提出的BCNN-IDS~\cite{BCNN-IDS}是该方向的代表性工作。该方法将贝叶斯卷积神经网络（Bayesian CNN）应用于NIDS，通过变分推断（Variational Inference）和蒙特卡洛Dropout（MC Dropout）估计权重后验分布，在NSL-KDD和UNSW-NB15数据集上实现了高准确率的同时，为每个预测提供了不确定性度量。BCNN-IDS利用30个贝叶斯CNN子模型对输入流量进行判别，取各子模型Softmax输出的中位数作为最终预测，并设置阈值$\lambda$筛除置信度不足的决策。在二分类检测中，随着阈值提高（如$\lambda=0.9$），BCNN-IDS达到99\%以上的检测率且误报率低至约0.35\%。该方法区分了两类不确定性：\textbf{认知不确定性}（epistemic uncertainty，源于模型对数据分布的学习不足，可通过增加训练数据或改进模型结构来降低）和\textbf{偶然不确定性}（aleatoric uncertainty，源于数据本身的固有噪声，无法通过增加数据消除）。这种区分对于入侵检测具有重要意义：高认知不确定性提示模型对该样本缺乏经验，需要人工审查或主动学习；高偶然不确定性则提示数据本身存在噪声或对抗性扰动。

然而，BCNN-IDS仍基于独立的流量特征向量进行分类，未显式建模流量间的关系结构或时序关联，难以捕获协同攻击或多步攻击链等复杂模式。这一局限在后续研究中逐渐被认识到。

\paragraph{（2）不确定性感知训练与校准}

为提高模型对不确定性的刻画能力，后续研究开始关注不确定性感知训练与校准。Tabarisaadi等人（2024）提出了不确定性感知训练框架~\cite{Tabarisaadi2024}，引入多目标损失函数同时优化预测精度和不确定性质量。该框架在传统交叉熵损失基础上增加一项KL散度惩罚，拉大模型对正确预测与错误预测的不确定性分布差异，使模型对正确分类更确信、对错误分类更犹疑。训练过程中，模型被鼓励对准确预测输出低熵，对发生误判的样本输出高熵，从而显著降低了过度自信的错误判断率。实验证明，与传统仅关注准确率的训练相比，该方法在多个数据集上提高``不确定性准确率(UA)''，即模型能更可靠地区分自身何时可信。

类似地，Li等人（2025）构建了可信贝叶斯深度学习框架（TMFD）~\cite{Li2025TrustworthyBayesian}，通过结合贝叶斯神经网络与置信度校准策略提升模型决策的可靠性。该框架采用贝叶斯CNN对传感信号分类，并利用$\alpha$散度变分推断获得参数后验。在此基础上，TMFD将预测不确定性细分为认知不确定性和偶然不确定性：通过蒙特卡洛Dropout采样计算预测分布熵以量化数据噪声不确定性，计算预测输出的互信息以量化模型认知不确定性，从而识别分布外样本（未知故障）。此外，TMFD引入了置信度校准机制，在训练中将校准误差纳入损失函数（AMCL损失），并在推断后结合温度缩放和等径回归进行后处理校准，大幅降低了预期校准误差ECE。通过在机械故障诊断数据上的实验，TMFD证明了提高模型置信度与实际准确率的一致性对于检测未知故障的重要意义。这些不确定性优化与校准的方法同样适用于网络安全领域，有望缓解深度模型对新型攻击过度自信的问题，提升对异常流量的检测能力。

\paragraph{（3）贝叶斯图神经网络的初步尝试}

与此同时，研究者也开始探索图结构建模与贝叶斯深度学习的结合，以更好地表示网络流量的复杂关系。传统IDS大多将每条流量记录视作独立样本，未利用主机、服务、时间等关联信息；而图神经网络（GNN）能够天然刻画网络拓扑和实体关系，在入侵检测中展现出潜力。然而，此类模型通常采用确定性学习，缺乏不确定性量化机制，且往往假设拓扑结构静态不变。

在贝叶斯图神经网络的理论研究方面，Zhang等人（2019）在AAAI会议上提出了贝叶斯图卷积网络框架~\cite{Zhang2019BGCN}，将观测图视为参数化随机图的实例，对图结构和GNN权重进行联合后验推断，证明了对图结构进行贝叶斯推断的可行性。Hasanzadeh等人（2020）在ICML提出Graph DropConnect方法~\cite{Hasanzadeh2020GDC}，利用自适应连接采样来近似Bayesian GNN训练，为每条边引入随机采样（服从Beta分布的DropEdge策略），从而达到正则化深层GNN、防止过拟合，并能够建模图结构不确定性。在推荐系统领域，Sun等人（华为诺亚方舟实验室）提出了贝叶斯图卷积框架（BGCF）~\cite{Sun2020}，针对用户-物品二分图存在的数据稀疏和噪声交互问题，将观测到的交互图视为不完备的随机实现，通过节点复制机制生成多个扰动图样本，这种对图结构不确定性的显式建模有效提升了推荐结果的准确性和多样性。

类似思想正逐步影响网络安全领域：一些最新工作尝试将通信网络抽象为图，并利用GNN捕获流量的时空依赖。然而，这些基于图的安全模型大多尚未结合贝叶斯不确定性估计。当前鲜有工作将图结构建模、时序动态和贝叶斯不确定性量化有机融合，以同时应对5G/6G异构网络中的复杂拓扑和未知攻击检测挑战。

\paragraph{（4）知识引导的入侵检测探索}

近期研究开始探索将领域知识融入入侵检测模型，以提升模型的鲁棒性和可迁移性。例如，K-GetNID（2024~\cite{Wang2024KGetNID}利用网络安全领域知识（如攻击阶段、资产角色等）指导图模型构建，通过知识图谱辅助特征学习，提升了早期检测能力和跨域泛化性能。该方法证明了将结构化领域知识融入图神经网络的有效性，但其主要关注知识引导的特征学习，并未引入贝叶斯不确定性量化机制，无法为预测提供可信度评估。类似地，2025年的GNN-IDS~\cite{sun2024gnnids}多视角入侵检测框架将协议、主机、攻击行为等关联构建为知识图谱，与深度学习提取的特征融合，提升了模型的空间和关系特征表达能力，但同样缺乏不确定性量化。

这些工作表明，将领域知识作为先验融入模型是提升入侵检测性能的有效途径，但如何将知识引导与贝叶斯不确定性量化有机结合，实现知识先验与不确定性估计的协同优化，仍是该领域的前沿挑战。

\paragraph{（5）现有工作的系统性不足}

综合来看，现有工作在贝叶斯深度学习赋能网络安全方面虽取得初步进展，但仍存在诸多共同不足：

\textbf{(a) 结构化建模不足}——许多入侵检测方法仅在样本层面应用贝叶斯神经网络或不确定性近似，忽略了主机之间、流量之间的拓扑关系和攻击的时序链条，未充分利用图结构信息。这使得模型难以捕获分布式、多阶段攻击的全局模式。

\textbf{(b) 时空依赖建模缺失}——现有贝叶斯NIDS工作主要聚焦于传统的流统计特征或图像化流量表示，输入仍为固定维度的流特征向量，无法捕捉网络流量的时空依赖和拓扑结构。即使部分工作引入了图结构，也往往假设静态拓扑，未能建模攻击的时间演化。

\textbf{(c) 先验设计不充分}——大部分贝叶斯深度模型采用默认简单先验（如各参数独立的高斯分布），缺乏对安全领域先验知识的融合。例如，不同协议字段、端口号等特征间的已知关联并未体现在模型先验中，导致模型后验难以及时校正过度自信的问题。虽然近期研究（如K-GetNID）开始探索知识引导的图模型，但这些工作并未将领域知识作为贝叶斯先验融入模型，而是通过知识图谱辅助特征学习，且缺乏不确定性量化机制。

\textbf{(d) 可解释性与不确定性脱节}——现有研究往往只给出不确定性数值或利用不确定性筛除低置信度输出，但很少结合不确定性信息去解释模型决策或揭示异常样本的成因。例如，对于预测置信度低或不确定性高的流量，没有机制追踪其涉及的关键特征关系，使安全分析人员难以依据模型提供的信息作出判断。

\textbf{(e) 计算效率挑战}——贝叶斯推断的计算开销（如多次前向传播采样）在实时检测场景下仍是挑战。现有方法大多采用固定次数的蒙特卡洛采样（如30次），未考虑自适应采样策略以平衡精度与效率。

上述不足限制了贝叶斯深度学习模型在开放世界网络安全场景中的实用性，也为本研究提供了明确的改进方向。


\subsubsection{本课题的技术空缺与创新点}

\paragraph{（1）核心技术空缺}


综合上述两方面的研究现状，当前网络安全领域存在以下核心技术空缺：

\textbf{空缺一：缺乏统一的包级时空图建模框架}——现有时空图神经网络方法大多依赖流级抽象或主机级拓扑，需要为每条流维护状态表和图结构，在高并发5G/6G场景下面临$O(M)$复杂度和内存开销瓶颈。如何在无需流状态维护的前提下，直接在包序列上构建时空图，同时捕捉攻击的时间演化（前后包的关联）和空间聚类（相似包的结构），是该领域的首要技术空缺。

\textbf{空缺二：缺乏贝叶斯时空图神经网络}——现有贝叶斯深度学习方法主要应用于传统的流统计特征或图像化流量表示，鲜有工作将其与时空图神经网络深度结合。虽然已有研究在推荐系统等领域探索了贝叶斯图神经网络（如BGCF~\cite{Sun2020}、Graph DropConnect~\cite{Hasanzadeh2020GDC}），但在网络安全场景中，如何在时空图神经网络中引入变分推断，使模型在保持高分类精度的同时，输出可信的不确定性度量（区分认知不确定性和偶然不确定性），同时兼顾实时检测场景下的计算效率，仍是该领域的关键技术空缺。

\textbf{空缺三：缺乏领域知识先验与不确定性量化的融合}——近期研究开始探索将领域知识融入入侵检测模型，例如K-GetNID（TIFS 2024）利用网络安全领域知识（如攻击阶段、资产角色等）指导图模型构建，提升了早期检测能力和跨域泛化性能。然而，这些工作主要关注知识引导的特征学习，并未引入贝叶斯不确定性量化机制。如何将结构化领域知识（如特征依赖图、攻击链知识图谱）作为先验融入贝叶斯时空图神经网络，同时实现不确定性量化与知识引导的协同优化，是该领域的前沿技术空缺。

\textbf{空缺四：缺乏不确定性驱动的安全决策机制}——现有研究往往只给出不确定性数值，但很少将不确定性转化为可操作的安全决策支持。如何将不确定性量化与告警优先级、人工复核策略、分布外（OOD）检测、主动学习等安全运维流程结合，是该领域的应用技术空缺。

\paragraph{（2）本课题的创新点}


针对上述技术空缺，本课题提出贝叶斯时空图卷积网络（BSTGCN）框架，具有以下创新点：

\textbf{创新一：无流状态维护的包级时空图建模}——本课题提出一种新颖的包级时空图建模范式：(1)采用Transformer编码器对包序列进行全局编码，无需为每条流维护独立状态表；(2)通过KNN图构建器（$k=10$）基于余弦相似度动态构建包间空间图，捕捉相似包的结构聚类；(3)采用贝叶斯时空卷积块（Bayesian ST Block）并行建模时间依赖（因果扩张卷积，膨胀率$2^i$）和空间依赖（GraphSAGE聚合），通过门控融合机制自适应权衡时空特征。这种设计在保持包级时间分辨率（$10^{-3}\!\sim\!10^{-2}$秒）的同时，避免了流表维护的$O(M)$开销，适合高并发6G网络场景。与已有流级图嵌入方法（如Graph2vec+RF~\cite{Graph2vec+RF}、GTAE-IDS~\cite{GTAE-IDS}）相比，本方法无需为每条流维护独立状态表，显著降低了内存开销和流表查找复杂度。

\textbf{创新二：端到端的贝叶斯时空图神经网络}——本课题系统性地将变分推断引入时空图神经网络，构建端到端的贝叶斯框架：(1)在输入投影层、时空卷积块、门控注意力池化层和分类器中引入贝叶斯线性层（BayesianLinear），将权重建模为高斯分布$\mathcal{N}(\mu, \sigma^2)$；(2)采用重参数化技巧（$W = \mu + \sigma \odot \epsilon$，$\epsilon \sim \mathcal{N}(0, I)$）实现可微采样；(3)设计ELBO损失函数$\mathcal{L} = \mathcal{L}_{\text{rec}} + \beta(t) \cdot \text{KL}[q(\Theta) \| p(\Theta)]$，其中$\beta(t) = \beta_{\max} \cdot \min(1, t/T_{\text{warmup}})$采用KL退火策略（$\beta_{\max}=5\times10^{-4}$，$T_{\text{warmup}}=60$）避免后验坍缩；(4)在推理阶段，针对实时检测需求，结合蒙特卡洛采样与置信度阈值设计自适应推理流程（初始采样$n_{\min}=3$次，若置信度$<0.95$则继续至$n_{\max}=8$次），在保证不确定性估计质量的同时减少平均采样次数约40\%和推理延迟。这种设计使模型能够输出认知不确定性（互信息$I(Y;W|X)$）和偶然不确定性（预测熵$H(Y|X)$），为零信任架构提供可信度评估，同时兼顾实时检测的效率要求。与已有贝叶斯入侵检测方法（如BCNN-IDS~\cite{BCNN-IDS}）相比，本方法将贝叶斯推断与时空图建模深度融合，能够同时捕捉流量的关系结构和时序演化，并量化预测不确定性。

\textbf{创新三：不确定性驱动的安全决策支持}——本课题将不确定性量化与安全运维流程深度结合：(1)告警优先级：根据预测概率和不确定性计算风险分数$\text{risk\_score} = 1 - \max(p) + \alpha \cdot I(Y;W|X)$，优先处理高风险告警；(2)人工复核策略：对不确定性超过阈值的样本触发人工审查，避免误报和漏报；(3)OOD检测：利用高认知不确定性识别分布外样本（未知攻击类型），触发模型更新或安全响应；(4)主动学习：选择高不确定性样本进行标注和再训练，持续提升模型性能。这种设计将不确定性从``模型内部指标''转化为``可操作的决策支持''，提升了模型的实用性和可信度。


\subsection{研究目标与拟解决的关键问题}

本课题的最终目标是：构建一个贝叶斯时空图卷积网络（BSTGCN）框架，实现高精度、多类入侵检测的同时，提供可解释且可度量的不确定性信息，为6G网络安全中的可信入侵检测提供系统性解决方案。

具体而言，本课题拟解决以下三个关键问题：

\subsubsection{关键问题一：包级时空图建模}

\textbf{问题描述}：如何在包级序列上构建时空图，同时捕捉攻击的时间演化（前后包的关联）和空间聚类（相似包的结构），且无需为每条流维护独立状态表？

\textbf{解决方案}：本课题提出一种新颖的包级时空图建模范式：

(1)采用Transformer编码器（$L=50$层，$D=25$维）对包序列进行全局编码，利用自注意力机制捕捉长距离依赖；

(2)通过KNN图构建器（$k=10$）基于余弦相似度动态构建包间空间图，捕捉相似包的结构聚类；

(3)采用贝叶斯深度可分离时间卷积网络（Bayesian Depthwise Separable TCN）建模时间依赖，采用贝叶斯GraphSAGE层建模空间依赖，通过门控融合机制自适应权衡时空特征。

这种设计在保持包级时间分辨率的同时，避免了流表维护的$O(M)$开销。

\textbf{预期成果}：在5G-NIDD数据集上，BSTGCN的多类F1分数达到95\%以上，相比确定性基线（Deterministic STGCN）提升2-3个百分点，证明贝叶斯时空图建模的有效性。

\subsubsection{关键问题二：贝叶斯推断与不确定性量化及其效率优化}


\textbf{问题描述}：如何在时空图神经网络中引入变分推断，使模型在保持高分类精度的同时，输出可信的不确定性度量（区分认知不确定性和偶然不确定性），并兼顾实时检测场景下的计算效率？

\textbf{解决方案}：本课题在输入投影层、时空卷积块、门控注意力池化层和分类器中引入贝叶斯线性层，将权重建模为高斯分布$\mathcal{N}(\mu, \sigma^2)$。采用重参数化技巧实现可微采样，设计ELBO损失函数$\mathcal{L} = \mathcal{L}_{\text{rec}} + \beta(t) \cdot \text{KL}[q(\Theta) \| p(\Theta)]$，其中$\beta(t)$采用KL退火策略（$\beta_{\max}=5\times10^{-4}$，$T_{\text{warmup}}=60$）避免后验坍缩。在推理阶段，通过蒙特卡洛采样估计预测分布，计算认知不确定性（互信息$I(Y;W|X) = H(\mathbb{E}[p(Y|X,W)]) - \mathbb{E}[H(p(Y|X,W))]$）和偶然不确定性（预测熵$H(Y|X)$）。

针对实时检测的效率需求，本课题设计了自适应采样策略：(1)初始阶段采用最小采样次数（$n_{\min}=3$）快速估计预测分布；(2)若最大预测概率超过置信度阈值（$\text{confidence\_threshold}=0.95$），则提前终止采样；(3)否则继续采样至最大次数（$n_{\max}=8$），以获得更精确的不确定性估计。此外，采用混合精度训练（FP16）和确定性Transformer编码器（仅在关键层引入贝叶斯推断）进一步降低计算开销。

\textbf{预期成果}：在5G-NIDD数据集上，BSTGCN的不确定性估计与预测错误率呈现显著正相关（Spearman相关系数$>0.6$），证明不确定性量化的有效性。对于高不确定性样本（前10\%），模型的错误率显著高于低不确定性样本（后10\%），验证了不确定性作为可信度指标的可靠性。自适应采样策略将平均采样次数降低约40\%（从固定8次降至平均4.8次），推理速度提升约1.6倍，同时不确定性估计精度损失小于5\%，将推理延迟降低至可接受范围（<100ms/batch），满足实时检测需求。

\subsubsection{关键问题三：不确定性驱动的安全决策}


\textbf{问题描述}：如何将不确定性量化转化为可操作的安全决策支持，包括告警优先级、人工复核策略、OOD检测等？

\textbf{解决方案}：本课题将不确定性量化与安全运维流程深度结合：(1)告警优先级：根据预测概率和不确定性计算风险分数$\text{risk\_score} = 1 - \max(p) + \alpha \cdot I(Y;W|X)$（$\alpha=0.5$），优先处理高风险告警；(2)人工复核策略：对不确定性超过阈值（$I(Y;W|X) > 0.3$）的样本触发人工审查，避免误报和漏报；(3)OOD检测：利用高认知不确定性识别分布外样本（未知攻击类型），设置OOD阈值（$I(Y;W|X) > 0.5$）触发模型更新或安全响应；(4)主动学习：选择高不确定性样本进行标注和再训练，持续提升模型性能。

\textbf{预期成果}：在5G-NIDD数据集上，基于不确定性的告警优先级策略使安全分析人员在处理前20\%告警时即可覆盖80\%的真实攻击，显著提升运维效率。OOD检测在识别未知攻击类型时的AUROC达到0.85以上，证明不确定性驱动的安全决策支持的有效性。主动学习策略在标注10\%高不确定性样本后，模型F1分数提升3-5个百分点，验证了不确定性在持续学习中的价值。



\section{研究方法与技术路线}

\subsection{BSTGCN 整体框架}

BSTGCN 的核心思想是：\textbf{在包级序列上构建时空图，通过贝叶斯推断量化不确定性}。整体架构如下：

\begin{enumerate}
  \item \textbf{输入}：包序列 $(B, L, D)$，其中 $B$ 为批量大小，$L=50$ 为序列长度（包数），$D=25$ 为包特征维度；
  \item \textbf{确定性编码}：Transformer 编码器学习包的语义表示；
  \item \textbf{确定性图构建}：KNN 图构建器发现相似包的结构关系；
  \item \textbf{贝叶斯时空建模}：3 个贝叶斯 ST 块，每个块包含：
  \begin{itemize}
    \item 贝叶斯 TCN：捕捉时间依赖 + 参数不确定性；
    \item 贝叶斯 GraphSAGE：捕捉空间依赖 + 参数不确定性；
    \item 门控融合：自适应平衡时间和空间信息；
  \end{itemize}
  \item \textbf{贝叶斯聚合}：门控注意力池化，加权聚合序列节点；
  \item \textbf{贝叶斯分类}：分类器输出类别概率 + 认知不确定性。
\end{enumerate}

\textbf{设计哲学}：
\begin{itemize}
  \item \textbf{确定性 + 贝叶斯混合}：Transformer 和 KNN 图构建保持确定性（计算成本高，贝叶斯化收益有限），将贝叶斯化集中在时空建模层（核心不确定性来源）；
  \item \textbf{时空并行 + 门控融合}：TCN 和 GraphSAGE 并行处理，通过门控机制自适应平衡，避免简单串联或加权；
  \item \textbf{端到端训练}：ELBO 损失统一优化分类性能和不确定性量化。
\end{itemize}

\subsection{时空图建模：捕捉攻击的双重依赖}

\subsubsection{KNN 图构建：发现相似包的空间结构}

给定 Transformer 编码后的包表示 $\mathbf{X} \in \mathbb{R}^{B \times L \times D}$，我们构建 KNN 图 $\mathbf{A} \in \mathbb{R}^{B \times L \times L}$：

\begin{equation}
\mathbf{A}_{ij} = \begin{cases}
1, & \text{if } j \in \text{TopK}(\text{sim}(\mathbf{x}_i, \mathbf{x}_j), k) \\
0, & \text{otherwise}
\end{cases}
\end{equation}

其中 $\text{sim}(\mathbf{x}_i, \mathbf{x}_j) = \frac{\mathbf{x}_i \cdot \mathbf{x}_j}{\|\mathbf{x}_i\| \|\mathbf{x}_j\|}$ 为余弦相似度，$k=10$ 为邻居数。邻接矩阵对称化：$\mathbf{A} \leftarrow (\mathbf{A} + \mathbf{A}^\top) / 2$。

\textbf{直觉}：相似的包（如相同协议、相似长度）在特征空间中接近，KNN 图捕捉这种空间聚类关系。

\subsubsection{贝叶斯 GraphSAGE：空间依赖 + 参数不确定性}

GraphSAGE 通过邻居聚合捕捉空间依赖。给定节点特征 $\mathbf{X}$ 和邻接矩阵 $\mathbf{A}$，GraphSAGE 层计算：

\begin{align}
\mathbf{H}_{\text{neigh}} &= \text{AGG}(\{\mathbf{x}_j : j \in \mathcal{N}(i)\}) = \frac{1}{|\mathcal{N}(i)|} \sum_{j \in \mathcal{N}(i)} \mathbf{x}_j \\
\mathbf{h}_i^{(l+1)} &= \sigma\left( \mathbf{W}_{\text{self}}^{(l)} \mathbf{h}_i^{(l)} + \mathbf{W}_{\text{neigh}}^{(l)} \mathbf{H}_{\text{neigh}} \right)
\end{align}

在 BSTGCN 中，我们将权重 $\mathbf{W}_{\text{self}}$ 和 $\mathbf{W}_{\text{neigh}}$ 替换为贝叶斯线性层，权重服从变分后验分布 $q(\mathbf{W}) = \mathcal{N}(\boldsymbol{\mu}_W, \text{diag}(\boldsymbol{\sigma}_W^2))$。前向传播时，通过重参数化技巧采样权重：

\begin{equation}
\mathbf{W} = \boldsymbol{\mu}_W + \boldsymbol{\sigma}_W \odot \boldsymbol{\epsilon}, \quad \boldsymbol{\epsilon} \sim \mathcal{N}(0, \mathbf{I})
\end{equation}

\textbf{直觉}：贝叶斯 GraphSAGE 不仅学习``如何聚合邻居''，还学习``对聚合权重的不确定性''，当训练数据不足或输入异常时，$\boldsymbol{\sigma}_W$ 增大，反映模型的不确定性。

\textbf{为什么用 GraphSAGE 而非 GATv2}：GraphSAGE 的邻居聚合是线性复杂度 $O(L \cdot k \cdot D)$，而 GATv2 的注意力机制是二次复杂度 $O(L^2 \cdot D)$。在包级图上（$L=50$），GATv2 的计算成本显著更高，而注意力机制的收益不明显（KNN 图已经选择了最相关的邻居）。因此，我们选择 GraphSAGE 提升效率，用门控融合实现自适应。

\subsubsection{贝叶斯 TCN：时间依赖 + 参数不确定性}

时间卷积网络（TCN）通过因果卷积和膨胀卷积捕捉长程时间依赖。我们采用深度可分离卷积降低参数量：

\begin{align}
\mathbf{H}_{\text{temp}} &= \text{DepthwiseConv1d}(\mathbf{X}) \\
\mathbf{H}_{\text{temp}} &= \text{PointwiseConv1d}(\mathbf{H}_{\text{temp}})
\end{align}

在 BSTGCN 中，卷积核权重同样替换为贝叶斯层，通过重参数化采样。膨胀率设置为 $\{1, 2, 4\}$（对应 3 个 ST 块），感受野覆盖 $\approx 20$ 个包。

\textbf{直觉}：贝叶斯 TCN 不仅学习``时间模式''，还学习``对时间模式的不确定性''，当包序列异常（如突发流量、罕见模式）时，模型输出高不确定性。

\subsubsection{门控融合：自适应平衡时空信息}

给定时间分支输出 $\mathbf{H}_t$ 和空间分支输出 $\mathbf{H}_s$，我们通过门控机制自适应融合：

\begin{align}
\mathbf{g}_t &= \sigma(\mathbf{W}_t \mathbf{X}), \quad \mathbf{g}_s = \sigma(\mathbf{W}_s \mathbf{X}) \\
\mathbf{g}_t &\leftarrow \frac{\mathbf{g}_t}{\mathbf{g}_t + \mathbf{g}_s + \epsilon}, \quad \mathbf{g}_s \leftarrow \frac{\mathbf{g}_s}{\mathbf{g}_t + \mathbf{g}_s + \epsilon} \\
\mathbf{H}^{(l+1)} &= \mathbf{g}_t \odot \mathbf{H}_t + \mathbf{g}_s \odot \mathbf{H}_s
\end{align}

其中 $\mathbf{W}_t$ 和 $\mathbf{W}_s$ 为贝叶斯线性层，$\sigma$ 为 Sigmoid 函数，$\epsilon=10^{-8}$ 为数值稳定项。

\textbf{直觉}：不同攻击类型对时间和空间信息的依赖不同（如端口扫描更依赖时序，DDoS 更依赖空间聚类），门控机制自适应调整权重。



\subsection{贝叶斯推断：量化认知不确定性}

\subsubsection{变分推断与 ELBO 损失}

给定数据集 $\mathcal{D} = \{(\mathbf{x}_i, y_i)\}_{i=1}^N$，贝叶斯推断的目标是学习参数的后验分布 $p(\Theta | \mathcal{D})$。根据贝叶斯定理：

\begin{equation}
p(\Theta | \mathcal{D}) = \frac{p(\mathcal{D} | \Theta) p(\Theta)}{p(\mathcal{D})} = \frac{p(\mathcal{D} | \Theta) p(\Theta)}{\int p(\mathcal{D} | \Theta) p(\Theta) d\Theta}
\end{equation}

由于分母（边缘似然）的积分难以计算，我们使用变分推断（Variational Inference），用变分后验 $q(\Theta)$ 近似真实后验 $p(\Theta | \mathcal{D})$。变分推断最大化证据下界（Evidence Lower Bound, ELBO）：

\begin{align}
\log p(\mathcal{D}) &\geq \mathbb{E}_{q(\Theta)}[\log p(\mathcal{D} | \Theta)] - \text{KL}[q(\Theta) \| p(\Theta)] \\
&= \mathcal{L}_{\text{ELBO}}
\end{align}

其中：
\begin{itemize}
  \item \textbf{第一项}：重建损失（Reconstruction Loss），衡量模型的分类性能；
  \item \textbf{第二项}：KL 散度（KL Divergence），衡量变分后验与先验的差异，起正则化作用。
\end{itemize}

\textbf{直觉}：ELBO 损失平衡两个目标：(1) 准确预测标签（重建损失），(2) 保持参数接近先验（KL 散度）。这使得模型在训练数据不足或输入异常时，后验方差增大，输出高不确定性。

在 BSTGCN 中，我们采用标签平滑交叉熵（Label Smoothing Cross-Entropy）作为重建损失：

\begin{equation}
\mathcal{L}_{\text{recon}} = -\frac{1}{N} \sum_{i=1}^N \sum_{c=1}^C \tilde{y}_{ic} \log p(y_i=c | \mathbf{x}_i, \Theta)
\end{equation}

其中 $\tilde{y}_{ic} = (1-\alpha) y_{ic} + \alpha / C$，$\alpha=0.1$ 为平滑系数，$C=9$ 为类别数。标签平滑防止模型过度自信，提升泛化能力。

KL 散度对所有贝叶斯层累加：

\begin{equation}
\text{KL}[q(\Theta) \| p(\Theta)] = \sum_{\text{layer}} \text{KL}[q(\mathbf{W}) \| p(\mathbf{W})]
\end{equation}

对于高斯先验 $p(\mathbf{W}) = \mathcal{N}(0, \sigma_{\text{prior}}^2 \mathbf{I})$ 和高斯变分后验 $q(\mathbf{W}) = \mathcal{N}(\boldsymbol{\mu}_W, \text{diag}(\boldsymbol{\sigma}_W^2))$，KL 散度有闭式解：

\begin{equation}
\text{KL} = \frac{1}{2} \sum_{i,j} \left[ \frac{\sigma_{W,ij}^2}{\sigma_{\text{prior}}^2} + \frac{\mu_{W,ij}^2}{\sigma_{\text{prior}}^2} - 1 - 2\log\frac{\sigma_{W,ij}}{\sigma_{\text{prior}}} \right]
\end{equation}

\textbf{直觉}：当 $\sigma_{W,ij}$ 远大于 $\sigma_{\text{prior}}$ 时，KL 散度增大，惩罚过大的后验方差；当 $\mu_{W,ij}$ 远离 0 时，KL 散度增大，惩罚偏离先验均值的参数。

\subsubsection{KL 权重退火：稳定训练}

直接优化 ELBO 会导致训练初期 KL 散度主导损失，模型收敛到先验附近，分类性能差（称为``后验塌陷''）。我们采用KL 权重退火（KL Weight Annealing）策略：

\begin{equation}
\mathcal{L} = \mathcal{L}_{\text{recon}} + \beta(t) \cdot \text{KL}, \quad \beta(t) = \beta_{\max} \cdot \min\left(1, \frac{t}{T_{\text{warmup}}}\right)
\end{equation}

其中 $\beta_{\max}=5\times10^{-4}$，$T_{\text{warmup}}=60$ 轮，$t$ 为当前训练轮数。

\textbf{训练过程}：
\begin{itemize}
  \item \textbf{初期}（$t < T_{\text{warmup}}$）：$\beta(t) \approx 0$，模型专注于分类，快速学习数据模式；
  \item \textbf{后期}（$t \geq T_{\text{warmup}}$）：$\beta(t) = \beta_{\max}$，模型平衡分类和正则化，学习不确定性。
\end{itemize}

\textbf{直觉}：KL 退火类似于``先学习，再反思``——先让模型学会分类，再让模型学会``对自己的预测有多确定``。

\subsubsection{蒙特卡洛采样：估计预测分布}

训练完成后，我们通过蒙特卡洛（MC）采样估计预测分布。给定输入 $\mathbf{x}$，我们从变分后验 $q(\Theta)$ 中采样 $M$ 次参数，得到 $M$ 个预测：

\begin{equation}
\Theta^{(m)} \sim q(\Theta), \quad p^{(m)} = p(y | \mathbf{x}, \Theta^{(m)}), \quad m=1,\ldots,M
\end{equation}

预测分布的均值和方差为：

\begin{align}
\bar{p}_c &= \frac{1}{M} \sum_{m=1}^M p_c^{(m)}, \quad c=1,\ldots,C \\
\text{Var}(p_c) &= \frac{1}{M} \sum_{m=1}^M (p_c^{(m)} - \bar{p}_c)^2
\end{align}

最终预测为 $\hat{y} = \arg\max_c \bar{p}_c$。

\textbf{直觉}：MC 采样通过多次``试探``，估计模型对输入的预测分布。如果 $M$ 次预测一致（方差小），说明模型确定；如果 $M$ 次预测分散（方差大），说明模型不确定。

\subsubsection{不确定性量化：互信息}

我们使用互信息（Mutual Information）量化认知不确定性（Epistemic Uncertainty），即模型参数的不确定性：

\begin{equation}
\mathcal{U}_{\text{epi}}(\mathbf{x}) = \mathcal{H}\left[\mathbb{E}_{\Theta}[p(y|\mathbf{x},\Theta)]\right] - \mathbb{E}_{\Theta}\left[\mathcal{H}[p(y|\mathbf{x},\Theta)]\right]
\end{equation}

其中 $\mathcal{H}[p] = -\sum_c p_c \log p_c$ 为熵。互信息分解为两项：
\begin{itemize}
  \item 第一项：预测分布均值的熵，衡量``平均预测``的不确定性；
  \item 第二项：预测分布熵的均值，衡量``每次预测``的平均不确定性。
\end{itemize}

\textbf{直觉}：
\begin{itemize}
  \item 如果 $M$ 次预测都是 $[0.9, 0.1]$（一致但不确定），第一项高，第二项高，互信息低；
  \item 如果 $M$ 次预测分别是 $[1, 0]$ 和 $[0, 1]$（分散但确定），第一项高，第二项低，互信息高；
  \item 互信息高说明模型``不知道该相信哪个预测``，即认知不确定性高。
\end{itemize}

在实现中，我们使用 PyTorch 的 \texttt{torch.special.entr} 函数计算熵，避免 $p=0$ 时的数值问题：

\begin{equation}
\mathcal{H}[p] = \sum_c \text{entr}(p_c), \quad \text{entr}(p) = \begin{cases}
-p \log p, & p > 0 \\
0, & p = 0
\end{cases}
\end{equation}

\subsubsection{自适应 MC 采样：平衡精度与效率}

标准 MC 推断需要固定采样次数（如 $M=30$），但不同样本的不确定性估计收敛速度不同。我们提出自适应 MC 采样策略：

% \begin{algorithm}[H]
% \caption{自适应 MC 采样}
% \begin{algorithmic}[1]
% \STATE 初始采样 $M_{\min}=3$ 次，计算 $\bar{p}$ 和 $\mathcal{U}_{\text{epi}}$
% \IF{$\max_c \bar{p}_c > \tau_{\text{conf}}$ \AND $\mathcal{U}_{\text{epi}} < \tau_{\text{unc}}$}
%     \STATE 停止采样（高置信度 + 低不确定性）
% \ELSE
%     \STATE 继续采样至 $M_{\max}=8$（验证）或 $M_{\max}=30$（测试）
% \ENDIF
% \STATE 返回 $\bar{p}$ 和 $\mathcal{U}_{\text{epi}}$
% \end{algorithmic}
% \end{algorithm}

其中 $\tau_{\text{conf}}=0.95$，$\tau_{\text{unc}}=0.05$。

\textbf{直觉}：对于``模型很确定``的样本（如正常流量），3 次采样足够；对于``模型不确定``的样本（如边界样本、未知攻击），需要更多采样以准确估计不确定性。

\textbf{效果}：在 5G-NIDD 数据集上，自适应采样将验证时间从 30 秒/epoch 减少到 12 秒/epoch，减少约 60\%，而分类性能几乎不变（<0.1\% 差异）。

\subsection{训练与评估}

\subsubsection{优化器配置}

我们使用 AdamW 优化器，配置如下：

\begin{itemize}
  \item 学习率：$\eta=5\times10^{-4}$
  \item 权重衰减：$\lambda=10^{-4}$
  \item 梯度裁剪：$\|\nabla\|_{\max}=0.5$（防止梯度爆炸）
\end{itemize}

学习率调度采用余弦退火 + 预热（Cosine Annealing with Warmup）：

\begin{equation}
\eta(t) = \begin{cases}
\eta_{\max} \cdot \frac{t}{T_{\text{warmup}}}, & t \leq T_{\text{warmup}} \\
\eta_{\min} + \frac{1}{2}(\eta_{\max} - \eta_{\min})\left(1 + \cos\frac{\pi(t-T_{\text{warmup}})}{T_{\max}-T_{\text{warmup}}}\right), & t > T_{\text{warmup}}
\end{cases}
\end{equation}

其中 $T_{\text{warmup}}=10$ 轮，$T_{\max}=150$ 轮，$\eta_{\min}=10^{-6}$。

\textbf{直觉}：预热阶段让模型稳定收敛，余弦退火阶段让模型精细调整。

\subsubsection{类别权重}

5G-NIDD 数据集存在类别不平衡（正常流量占 60\%，某些攻击类型仅占 1\%）。我们采用逆频率加权（Inverse Frequency Weighting）：

\begin{equation}
w_c = \frac{N}{C \cdot N_c}, \quad c=1,\ldots,C
\end{equation}

其中 $N$ 为总样本数，$N_c$ 为类别 $c$ 的样本数，$C=9$ 为类别数。对于极少数类别（如 Scanning），我们额外增强权重：

\begin{equation}
w_c \leftarrow w_c \cdot \left(1 + \alpha \cdot \mathbb{I}[N_c < \tau]\right)
\end{equation}

其中 $\alpha=0.5$，$\tau=5000$。

\textbf{直觉}：逆频率加权让模型更关注少数类别，避免模型偏向多数类别。



\subsubsection{评估指标}

我们从三个维度评估 BSTGCN：

\textbf{1. 分类性能}：
\begin{itemize}
  \item 准确率（Accuracy）：$\text{Acc} = \frac{\text{正确预测数}}{\text{总样本数}}$
  \item Macro F1：各类别 F1 分数的算术平均，衡量对少数类别的识别能力
  \item Weighted F1：各类别 F1 分数的加权平均，权重为类别频率
  \item Per-class F1：每个攻击类型的 F1 分数，识别模型的弱点
\end{itemize}

\textbf{2. 不确定性量化}：
\begin{itemize}
  \item 不确定性-错误相关性：Pearson 相关系数 $\rho(\mathcal{U}_{\text{epi}}, \mathbb{I}_{\text{error}})$，衡量不确定性与错误的相关性。理想情况下，错误样本应有更高的不确定性（$\rho > 0$）
  \item 不确定性分布：正确样本 vs 错误样本的不确定性分布，通过 KS 检验判断两者是否显著不同
  \item 高不确定性样本分析：Top-K 高不确定性样本的错误率，验证不确定性的实用价值
\end{itemize}

\textbf{3. 计算效率}：
\begin{itemize}
  \item 训练时间：每轮训练时间（秒/epoch）
  \item 验证时间：每轮验证时间（秒/epoch）
  \item 推断延迟：单样本推断时间（毫秒/样本）
  \item 模型参数量：总参数数量，衡量模型复杂度
\end{itemize}

\section{实验验证与结果分析}

\subsection{数据集与实验设置}

\subsubsection{5G-NIDD 数据集}

5G-NIDD（5G Network Intrusion Detection Dataset）是一个面向 5G 网络的入侵检测数据集，包含真实5G基站环境下采集的流量数据。数据集共包含 1,215,890 条流记录（来自两个基站BS1和BS2），每条流记录包含 112 个特征维度，涵盖网络五元组、流量统计、时间特征和5G特有的QoS标记等信息。

\textbf{攻击类型分布：}数据集包含 1 类良性流量和 8 类攻击流量，具体分布如下：

\begin{itemize}
  \item \textbf{良性流量}（Benign）：477,737 条记录（39.2\%）
  \item \textbf{UDPFlood}：457,340 条记录（37.5\%）
  \item \textbf{HTTPFlood}：140,812 条记录（11.5\%）
  \item \textbf{SlowrateDoS}：73,124 条记录（6.0\%）
  \item \textbf{TCPConnectScan}：20,052 条记录（1.6\%）
  \item \textbf{SYNScan}：20,043 条记录（1.6\%）
  \item \textbf{UDPScan}：15,906 条记录（1.3\%）
  \item \textbf{SYNFlood}：9,721 条记录（0.8\%）
  \item \textbf{ICMPFlood}：1,155 条记录（0.1\%）
\end{itemize}

\textbf{5G网络特征：}数据集包含5G网络特有的语义信息，包括：(1) SCTP协议流量（占0.53\%），用于5G信令传输；(2) QoS标记字段（sTos/dTos），包含cs0、ef、af11等服务质量标识；(3) 真实的5G网络拓扑结构，包括UE网段（10.155.15.x）、基站核心网（10.41.150.x）和外部服务IP的三层架构。

\textbf{时间特征：}数据集中99.6\%的流持续时间超过4.5秒，属于长流，适合研究攻击在时间维度上的演化模式。协议分布为TCP（89.15\%）、UDP（6.14\%）、ICMP（4.16\%）、SCTP（0.53\%）和ARP（0.01\%）。

我们从原始流记录中提取包级序列，每条流选取连续 50 个数据包，每个数据包提取 25 维特征（协议、端口、方向、长度、时间间隔等）。数据集划分如下：

\begin{table}[h]
\centering
\begin{tabular}{lrrr}
\hline
\textbf{划分} & \textbf{样本数} & \textbf{比例} & \textbf{用途} \\
\hline
训练集 & 150,154 & 70\% & 模型训练 \\
验证集 & 32,176 & 15\% & 超参数调优、早停 \\
测试集 & 32,177 & 15\% & 最终性能评估 \\
\hline
\textbf{总计} & \textbf{214,507} & \textbf{100\%} & \\
\hline
\end{tabular}
\caption{5G-NIDD 数据集划分}
\end{table}

\subsubsection{5G-NIDD 相较于传统数据集的优势}

相较于文献综述中提到的传统入侵检测数据集（UNSW-NB15、BoT-IoT、CSE-CIC-IDS2017、CSE-CIC-IDS2018），5G-NIDD 数据集在以下四个方面更适合本研究的时空图神经网络建模：

\textbf{（1）真实5G网络环境：}传统数据集（如CIC-IDS系列）主要采集自4G或传统互联网环境，缺乏5G特有的网络特征。5G-NIDD包含真实5G基站（BS1、BS2）采集的流量，保留了SCTP信令协议、QoS标记（sTos/dTos字段）、UE-gNB-Core三层网络拓扑等5G语义信息。这些特征对于研究5G网络中的攻击传播模式和控制面/用户面的异常行为至关重要，是传统数据集无法提供的。

\textbf{（2）长流占比高，适合时序建模：}5G-NIDD中99.6\%的流持续时间超过4.5秒，属于长流，包含丰富的时间演化信息。相比之下，UNSW-NB15和BoT-IoT等数据集中短流占比较高，流内包序列较短，难以充分验证时空图神经网络对长时间依赖的建模能力。长流特性使得本研究提出的TCN-Transformer时序编码器能够捕捉攻击在多个时间尺度上的演化模式（如SlowrateDoS的慢速攻击特征）。

\textbf{（3）攻击类型覆盖5G典型威胁：}5G-NIDD包含的8类攻击（UDPFlood、HTTPFlood、SlowrateDoS、TCPConnectScan、SYNScan、UDPScan、SYNFlood、ICMPFlood）覆盖了5G网络面临的典型威胁场景，包括针对用户面的洪泛攻击（UDPFlood、HTTPFlood）、针对控制面的扫描攻击（SYNScan、UDPScan）以及慢速攻击（SlowrateDoS）。这些攻击类型在时空维度上呈现不同的模式：洪泛攻击表现为短时间内大量包的突发，扫描攻击表现为跨多个目标IP的空间扩散，慢速攻击表现为长时间的低速率异常。这种多样性为验证BSTGCN的时空建模能力提供了理想的测试场景。

\textbf{（4）网络拓扑结构清晰，支持空间图构建：}5G-NIDD的IP地址分布呈现清晰的层次结构（UE网段10.155.15.x、基站核心网10.41.150.x、外部服务IP），使得基于KNN的动态图构建能够捕捉真实的网络拓扑关系。传统数据集（如CIC-IDS系列）往往缺乏明确的网络拓扑信息，IP地址分布较为随机，难以验证空间图神经网络对攻击传播路径的建模效果。此外，5G-NIDD包含多个攻击工具（Hping3、Goldeneye、Torshammer、Nmap、Slowloris）生成的流量，攻击源和目标的空间分布更加多样化，有助于评估模型对分布式攻击的检测能力。

综上，5G-NIDD数据集的5G网络语义、长流特性、多样化攻击类型和清晰的网络拓扑结构，使其成为验证本研究提出的贝叶斯时空图卷积网络（BSTGCN）的理想数据集。相较于传统数据集，5G-NIDD能够更全面地评估模型在真实5G环境下对时空依赖、攻击演化和不确定性量化的建模能力。

\subsubsection{5G数据集对6G研究的适用性论证}

尽管本研究面向6G网络安全，但选择5G-NIDD数据集作为实验验证基础具有充分的合理性和必要性，主要基于以下四个方面的考虑：

\textbf{（1）6G网络的演进继承性：}根据ITU-R IMT-2030愿景~\cite{ITU_IMT2030_2023}和IMT-2030(6G)推进组白皮书~\cite{IMT2030_6G_WhitePaper_2021}，6G网络并非对5G的完全颠覆，而是在5G基础上的演进增强。6G将继承5G的核心网络架构（如服务化架构SBA、网络切片、边缘计算等），并在此基础上引入新的技术特性（如太赫兹通信、空天地一体化、通感算融合等）。从网络安全角度看，5G面临的主要威胁类型（洪泛攻击、扫描攻击、慢速攻击等）在6G中不仅会持续存在，还可能因网络复杂度提升而更加严峻。因此，在5G数据集上验证的入侵检测方法，其核心建模能力（时空依赖建模、不确定性量化）可以直接迁移到6G场景。

\textbf{（2）6G真实数据集的缺失：}截至目前（2025年），6G网络仍处于标准制定和技术研发阶段，全球尚未部署商用6G网络，因此不存在真实的6G网络流量数据集。现有的6G相关研究主要依赖仿真数据或理论分析，但仿真数据往往无法完全反映真实网络环境的复杂性和攻击行为的多样性。相比之下，5G-NIDD数据集来自真实5G基站环境，包含真实的网络拓扑、协议栈和攻击流量，能够为模型提供高质量的训练和验证数据。在6G商用数据集可用之前，使用5G数据集是当前学术界和工业界的普遍做法，也是推动6G安全技术预研的必要选择。

\textbf{（3）方法论的通用性与可迁移性：}本研究的核心贡献在于提出一种通用的贝叶斯时空图神经网络框架（BSTGCN），而非针对特定网络代际的定制化方法。BSTGCN的核心建模能力——包级时空图构建、变分推断、不确定性量化——是与具体网络技术（5G或6G）解耦的。模型的输入是包级特征序列（协议、端口、长度、时间间隔等），这些特征在5G和6G中具有相同的语义。因此，在5G数据集上训练的BSTGCN模型，可以通过迁移学习或增量学习的方式，快速适配到6G网络环境。本研究在5G数据集上验证方法的有效性，为未来在6G数据集上的应用奠定了坚实的方法论基础。

\textbf{（4）面向6G的前瞻性设计：}虽然实验数据来自5G网络，但BSTGCN的设计充分考虑了6G网络的特性和需求。首先，BSTGCN采用无流状态维护的包级建模范式，避免了流表维护的$O(M)$开销，能够适应6G网络的超高并发场景（每平方公里百万级设备连接）。其次，BSTGCN通过贝叶斯推断量化预测不确定性，满足6G零信任架构对可信AI的要求（持续验证、风险量化、可解释决策）。最后，BSTGCN的自适应采样策略将推理延迟控制在100ms以内，能够满足6G网络的超低时延要求（空口时延<1ms）。这些设计使得BSTGCN不仅适用于当前的5G网络，更具备向6G演进的技术储备。

\textbf{替代数据集的局限性分析：}除5G-NIDD外，现有的网络入侵检测数据集主要包括UNSW-NB15、BoT-IoT、CIC-IDS2017/2018等，但这些数据集均采集自4G或传统互联网环境，缺乏5G特有的网络特征（如SCTP协议、QoS标记、UE-gNB-Core三层拓扑）。使用这些传统数据集进行6G研究，其代际差距更大，技术相关性更低。相比之下，5G-NIDD作为最接近6G的真实网络数据集，能够更好地反映未来移动网络的流量特性和攻击模式。此外，5G-NIDD的长流特性（99.6\%流持续时间>4.5秒）和清晰的网络拓扑结构，为验证时空图神经网络的建模能力提供了理想的测试场景，这是传统数据集无法提供的。

综上所述，在6G真实数据集缺失的现状下，选择5G-NIDD数据集进行方法验证是合理且必要的。本研究通过在5G数据集上验证BSTGCN的有效性，为未来在6G网络中的应用提供了方法论基础和技术储备。随着6G网络的逐步部署和真实数据集的出现，BSTGCN可以通过迁移学习或增量学习的方式，快速适配到6G环境，实现从5G到6G的平滑演进。

\subsubsection{模型配置}

BSTGCN 的超参数配置如下：

\begin{table}[h]
\centering
\begin{tabular}{ll}
\hline
\textbf{超参数} & \textbf{值} \\
\hline
输入维度 & 25 \\
隐藏维度 & 128 \\
Transformer 头数 & 8 \\
ST 块数量 & 3 \\
TCN 膨胀率 & $\{1, 2, 4\}$ \\
KNN 邻居数 & 10 \\
GraphSAGE 聚合方式 & mean \\
先验标准差 & 0.5 \\
KL 权重最大值 & $5\times10^{-4}$ \\
KL 退火轮数 & 60 \\
Dropout & 0.15 \\
\hline
\textbf{总参数量} & \textbf{755,220} \\
\hline
\end{tabular}
\caption{BSTGCN 模型配置}
\end{table}

\subsubsection{训练配置}

\begin{itemize}
  \item \textbf{批量大小}：512
  \item \textbf{最大轮数}：150
  \item \textbf{早停}：验证集 Macro F1 连续 20 轮不提升则停止
  \item \textbf{优化器}：AdamW（$\eta=5\times10^{-4}$，$\lambda=10^{-4}$）
  \item \textbf{学习率调度}：余弦退火 + 10 轮预热
  \item \textbf{梯度裁剪}：$\|\nabla\|_{\max}=0.5$
  \item \textbf{标签平滑}：$\alpha=0.1$
  \item \textbf{类别权重}：逆频率加权
  \item \textbf{MC 采样}：训练时 1 次，验证时 3-8 次（自适应），测试时 30 次
\end{itemize}

\subsection{分类性能}

\subsubsection{整体性能}


\subsubsection{各类别性能}



\subsubsection{混淆矩阵}


\subsection{不确定性分析}

\subsubsection{不确定性-错误相关性}




\subsubsection{不确定性分布}


\subsubsection{高不确定性样本分析}



\subsection{计算效率}

\subsubsection{训练与推断时间}


\subsubsection{模型参数分布}

BSTGCN 的参数分布如下：

\begin{table}[h]
\centering
\begin{tabular}{lrr}
\hline
\textbf{模块} & \textbf{参数量} & \textbf{占比} \\
\hline
输入投影 & 3,328 & 0.4\% \\
Transformer（确定性） & 263,168 & 34.8\% \\
KNN 图构建（确定性） & 0 & 0.0\% \\
贝叶斯 ST 块 ×3 & 393,216 & 52.1\% \\
贝叶斯池化 & 33,024 & 4.4\% \\
贝叶斯分类器 & 62,484 & 8.3\% \\
\hline
\textbf{总计} & \textbf{755,220} & \textbf{100\%} \\
\hline
\end{tabular}
\caption{BSTGCN 参数分布}
\end{table}

\textbf{观察}：
\begin{itemize}
  \item 贝叶斯 ST 块占参数量的 52.1\%，是模型的核心；
  \item Transformer 占 34.8\%，但保持确定性（计算成本高，贝叶斯化收益有限）；
  \item 贝叶斯层（ST 块 + 池化 + 分类器）占 64.8\%，体现了``确定性 + 贝叶斯混合``的设计哲学。
\end{itemize}

\section{后续工作与研究计划}

\subsection{不确定性优化}



\subsection{OOD 检测与分布漂移}


\subsection{领域先验融合（FDG 引导）}


\textbf{注}：FDG 引导的先验设置是未来扩展方向，当前 BSTGCN 的核心贡献在于贝叶斯时空图建模框架。

\subsection{可解释性与安全决策}




\section{工作量与预期成果}

\subsection{已完成工作}

\subsubsection{理论研究（2024.09 - 2025.06）}

\begin{itemize}
  \item \textbf{文献调研}：系统调研时空图神经网络、贝叶斯深度学习、网络入侵检测领域的 100+ 篇论文；
  \item \textbf{理论框架}：设计 BSTGCN 的贝叶斯时空图建模框架，包括变分推断、MC 采样、不确定性量化；
  \item \textbf{技术选型}：确定 GraphSAGE（空间）+ TCN（时间）+ 门控融合的架构，平衡性能与效率。
\end{itemize}

\subsubsection{工程实现（2025.07 - 2025.11）}

\begin{itemize}
  \item \textbf{模型实现}：完成 BSTGCN 的 PyTorch 实现（约 2000 行代码），包括贝叶斯组件、KNN 图构建、自适应采样；
  \item \textbf{训练优化}：解决训练稳定性问题（KL 退火、梯度裁剪、标签平滑），将训练时间从 91 分钟优化至 60 分钟；
  \item \textbf{推断优化}：设计自适应 MC 采样策略，将验证时间减少 60\%（从 30 秒降至 12 秒/epoch）；
  \item \textbf{代码开源}：整理代码库，准备开源发布（GitHub）。
\end{itemize}

\subsubsection{实验验证（2025.11 - 2025.12）}

\begin{itemize}
  \item \textbf{数据集准备}：处理 5G-NIDD 数据集，提取包级序列（214,507 条流，50 包/流，25 特征/包）；
  \item \textbf{性能评估}：在测试集上达到 96.87\% 准确率、92.30\% Macro F1；
  \item \textbf{不确定性分析}：验证不确定性-错误相关性（$\rho=0.1156$），分析高不确定性样本；
  \item \textbf{效率分析}：量化自适应采样的效果（减少 60\% 验证时间，准确率仅下降 0.03\%）。
\end{itemize}

\subsection{后续工作计划}

\subsubsection{第一阶段：不确定性优化（2026.01 - 2026.03）}

\textbf{目标}：提升不确定性-错误相关性至 $\rho > 0.3$。

\textbf{任务}：
\begin{itemize}
  \item 实现多度量融合（互信息 + 预测熵 + 方差）；
  \item 实现不确定性校准（温度缩放、Platt Scaling）；
  \item 引入对抗训练，增强边界样本的不确定性感知；
  \item 在 5G-NIDD 和其他数据集（如 CIC-IDS2017）上验证。
\end{itemize}


\subsubsection{第二阶段：OOD 检测与分布漂移（2026.04 - 2026.06）}

\textbf{目标}：OOD 检测 AUROC > 0.85，分布漂移下准确率下降 < 5\%。

\textbf{任务}：
\begin{itemize}
  \item 按时间顺序划分数据集，模拟真实部署场景；
  \item 实现 OOD 检测算法（基于不确定性阈值、Mahalanobis 距离）；
  \item 引入概念漂移（Concept Drift），评估模型鲁棒性；
  \item 设计在线学习策略，适应分布漂移。
\end{itemize}



\subsubsection{第三阶段：领域先验融合（2026.07 - 2026.09）}

\textbf{目标}：利用 FDG 引导贝叶斯先验设置，提升可解释性和稳健性。

\textbf{任务}：
\begin{itemize}
  \item 将 FDG 的特征关系映射为贝叶斯先验的方差约束；
  \item 实现 FDG 引导的先验设置算法；
  \item 分析 FDG 引导对模型性能和可解释性的影响；
  \item 对比无先验、均匀先验、FDG 引导先验的效果。
\end{itemize}


\subsubsection{第四阶段：可解释性与安全决策（2026.10 - 2026.12）}

\textbf{目标}：将不确定性量化转化为可操作的安全决策工具。

\textbf{任务}：
\begin{itemize}
  \item 绘制 Risk-Coverage 曲线，量化``拒绝预测``策略的效果；
  \item 设计告警优先级策略（结合预测类别和不确定性）；
  \item 对高不确定性样本进行深度案例分析；
  \item 开发可视化工具，辅助安全运维人员决策。
\end{itemize}

\subsubsection{第五阶段：论文撰写与答辩（2027.01 - 2027.06）}

\textbf{目标}：完成博士论文撰写和答辩。

\textbf{任务}：
\begin{itemize}
  \item 整合前期工作，撰写博士论文（预计 150-200 页）；
  \item 准备答辩材料（PPT、演讲稿）；
  \item 参加预答辩，根据反馈修改论文；
  \item 正式答辩。
\end{itemize}



\subsection{预期成果}

\subsubsection{学术成果}

\begin{itemize}
  \item 期刊论文一篇：BSTGCN 核心框架（贝叶斯时空图建模 + 自适应采样），投稿 IEEE 系列期刊。
  \item 会议论文一篇，投稿通信领域顶会（IEEE ICC 或 Globecom）。
  \item 申请相关专利一项。
\end{itemize}



\subsubsection{应用价值}

\begin{itemize}
  \item \textbf{5G/6G 安全}：为 5G/6G 网络提供可信的入侵检测解决方案，支持零信任架构；
  \item \textbf{安全运维}：通过不确定性量化，辅助安全运维人员优先复核高风险样本，提升检测可靠性；
  \item \textbf{未知威胁检测}：通过 OOD 检测，识别未知攻击和分布漂移，增强模型的泛化能力；
  \item \textbf{可解释 AI}：通过 FDG 引导和可视化工具，提升模型的可解释性，增强用户信任。
\end{itemize}

\subsection{工作量总结}

\begin{table}[h]
\centering
\begin{tabular}{lll}
\hline
\textbf{阶段} & \textbf{时间} & \textbf{主要任务} \\
\hline
理论研究 & 2025.01 - 2025.06 & 文献调研、理论框架、技术选型 \\
工程实现 & 2025.07 - 2025.11 & 模型实现、训练优化、推断优化 \\
实验验证 & 2025.11 - 2025.12 & 数据集准备、性能评估、效率分析 \\
不确定性优化 & 2026.01 - 2026.02 & 多度量融合、校准、对抗训练 \\
OOD 检测 & 2026.03 - 2025.04 & 时间划分、OOD 检测、分布漂移 \\
领域先验融合 & 2026.05 - 2026.06 & FDG 引导先验设置、对比实验 \\
可解释性 & 2026.7- 2026.8 & Risk-Coverage、告警优先级、可视化 \\
论文撰写 & 2026.9 - 2026.12 & 博士论文撰写、预答辩、正式答辩 \\
\hline
\end{tabular}
\caption{工作量与时间规划}
\end{table}

\section{总结}

本课题提出 BSTGCN（Bayesian Spatio-Temporal Graph Convolutional Network），一个面向 6G 网络安全的贝叶斯时空图神经网络框架，聚焦于不确定性量化与可信入侵检测。BSTGCN 的核心贡献在于：

\begin{enumerate}
  \item \textbf{时空图建模}：通过 TCN + GraphSAGE + 门控融合，同时捕捉攻击行为的时间依赖和空间聚类；
  \item \textbf{贝叶斯推断}：通过变分推断 + MC 采样，量化认知不确定性，识别``模型不确定''的样本；
  \item \textbf{自适应优化}：通过自适应 MC 采样，平衡不确定性估计的精度与效率，减少 60\% 验证时间。
\end{enumerate}

在 5G-NIDD 数据集上，BSTGCN 达到 96.67\% 准确率、91.90\% Macro F1，不确定性-错误相关性为 0.1777。实验验证了 BSTGCN 的有效性：错误样本的平均不确定性比正确样本高 50.5\%，Top-1\% 高不确定性样本的错误率是全部样本的 6.3 倍。

后续工作将聚焦于不确定性优化、OOD 检测、领域先验融合（FDG 引导）、可解释性与安全决策，最终形成一套完整的可信 AI 入侵检测解决方案，为 5G/6G 网络的零信任架构提供技术支撑。


